\documentclass[11pt,a4paper]{article}
\usepackage[utf8]{inputenc}
\usepackage[T1]{fontenc}
\usepackage[margin=2cm]{geometry}
\usepackage{listings}
\usepackage{xcolor}
\usepackage{titlesec}
\usepackage{booktabs}
\usepackage{array}
\usepackage{enumitem}
\usepackage{fancyhdr}
\usepackage{mdframed}
\usepackage{amsmath}

% ── Colores ──────────────────────────────────────────────────────────────────
\definecolor{cppbg}{RGB}{245,245,245}
\definecolor{kw}{RGB}{0,0,180}
\definecolor{cm}{RGB}{80,120,80}
\definecolor{str}{RGB}{160,0,0}
\definecolor{rubgreen}{RGB}{0,120,60}
\definecolor{ruborange}{RGB}{200,80,0}

% ── Estilo de código ─────────────────────────────────────────────────────────
\lstset{
  backgroundcolor=\color{cppbg},
  basicstyle=\ttfamily\footnotesize,
  keywordstyle=\color{kw}\bfseries,
  commentstyle=\color{cm}\itshape,
  stringstyle=\color{str},
  breaklines=true,
  frame=single,
  rulecolor=\color{gray!40},
  numbers=left,
  numberstyle=\tiny\color{gray},
  numbersep=6pt,
  tabsize=4,
  showstringspaces=false,
  language=C++
}

% ── Cabecera ─────────────────────────────────────────────────────────────────
\pagestyle{fancy}
\fancyhf{}
\lhead{\small\textbf{Proyecto Integrador — Visión Estereoscópica}}
\rhead{\small Michael Lata · Jordy Romero — G8}
\cfoot{\small\thepage}

% ── Cajas de criterio ─────────────────────────────────────────────────────────
\newmdenv[linecolor=rubgreen,linewidth=1.5pt,
          frametitle={\color{rubgreen}\bfseries},
          frametitlebelowskip=4pt,
          innertopmargin=8pt]{critbox}

% ── Título ───────────────────────────────────────────────────────────────────
\title{
  {\Large\bfseries GUÍA DE DEFENSA — PROYECTO INTEGRADOR}\\[4pt]
  {\large Sistema de Visión Estereoscópica con ESP32-CAM}\\[4pt]
  {\normalsize Universidad Politécnica Salesiana · Visión Artificial G8}
}
\author{Michael Lata \and Jordy Romero}
\date{Junio 2026}

% =============================================================================
\begin{document}
\maketitle
\thispagestyle{fancy}

\begin{center}
\begin{tabular}{lcc}
\toprule
\textbf{Criterio} & \textbf{Ponderación} & \textbf{Archivo principal} \\
\midrule
1. Arquitectura de Captura y Hardware & 20\% & \texttt{ESP32\_CAM\_Stereo.ino, CameraSource.cpp} \\
2. Calidad del Mapa de Disparidad     & 25\% & \texttt{StereoProcessor.cpp} \\
3. Precisión y Estabilidad de la Distancia & 20\% & \texttt{KalmanDepth.cpp, StereoProcessor.cpp} \\
4. Efecto de Realidad Aumentada (AR) & 15\% & \texttt{ARRenderer.cpp} \\
5. Mini-Evaluación Individual (Defensa) & 20\% & (preguntas clave) \\
\midrule
\textbf{Total} & \textbf{100\%} & \\
\bottomrule
\end{tabular}
\end{center}

\vspace{6pt}
\noindent\textbf{Nota:} la ponderación del 85\% aplica con cámaras ESP32-CAM-MB/XIAO-SP; este proyecto usa ESP32-CAM AI-Thinker (\textbf{20\%}).

% =============================================================================
\section{Criterio 1 — Arquitectura de Captura y Hardware (20\%)}
% =============================================================================

\subsection*{Lo que pide la rúbrica}
Sistema desarrollado en OpenCV/C++, ESP32-CAM con parámetros de firmware ajustados
(AEC, AGC, AWB, XCLK), lectura concurrente de ambas cámaras con hilos para reducir latencia.

\subsection*{Firmware ESP32-CAM — \texttt{ESP32\_CAM\_Stereo.ino}}

\begin{critbox}[frametitle={Por qué AEC/AGC/AWB = OFF}]
SGBM calcula disparidad comparando la intensidad del mismo punto físico en el frame
izquierdo y derecho. Si cada cámara ajusta su exposición de forma independiente,
ese punto produce valores distintos en L y R → errores sistemáticos en el mapa.
Con exposición manual \textbf{igual en ambas cámaras} se garantiza paridad lumínica.
\end{critbox}

\begin{lstlisting}[caption={Parámetros fijos del firmware (líneas 52–58)}]
constexpr int  CAMERA_XCLK_HZ = 6000000; // 6 MHz - reduce jitter WiFi
constexpr bool CAMERA_AEC     = false;   // AUTO-EXPOSICION OFF -> manual
constexpr bool CAMERA_AGC     = false;   // AUTO-GANANCIA  OFF -> manual
constexpr bool CAMERA_AWB     = false;   // AUTO-BALANCE   OFF -> color fijo
constexpr int  CAMERA_EXPOSURE = 400;    // 0-1200; ajustar por ambiente
constexpr int  CAMERA_GAIN     = 0;      // 0-30; minimo = menos ruido
\end{lstlisting}

\begin{critbox}[frametitle={Por qué XCLK = 6 MHz (no 20 MHz)}]
A mayor frecuencia el sensor demanda más corriente y genera más calor, aumentando
el jitter en la transmisión WiFi. A 6 MHz el ESP32 prioriza la integridad del paquete,
reduciendo el desincronismo entre frames L y R que degrada la triangulación.
\end{critbox}

Los parámetros se pueden ajustar \textbf{sin reflashear} vía HTTP:
\texttt{GET /setexposure?v=400} y \texttt{GET /setgain?v=0}.

\subsection*{Hilos concurrentes — \texttt{CameraSource.cpp}}

\begin{lstlisting}[caption={Hilo productor de frames (HttpCameraSource)}]
// El constructor lanza el hilo worker inmediatamente
HttpCameraSource::HttpCameraSource(std::string url, int timeoutMs)
    : url_(url), timeoutMs_(timeoutMs) {
    worker_ = std::thread(&HttpCameraSource::workerLoop, this);
}
// El hilo descarga JPEG de /capture en bucle continuo
void HttpCameraSource::workerLoop() {
    while (running_.load()) {
        // ... curl descarga JPEG ...
        {
            std::lock_guard<std::mutex> lk(mutex_);
            latestFrame_ = decoded;  // protected write
            hasFrame_    = true;
        }
        framesReceived_.fetch_add(1);
    }
}
// El bucle principal lee sin bloquearse
bool HttpCameraSource::read(cv::Mat& out) {
    std::lock_guard<std::mutex> lk(mutex_);
    if (!hasFrame_) return false;
    out = latestFrame_.clone();
    return true;
}
\end{lstlisting}

\begin{itemize}[leftmargin=*]
  \item Cámara izquierda y cámara derecha tienen cada una su propio hilo productor.
  \item El hilo principal (controlador) consume los frames más recientes sin esperar.
  \item El \texttt{mutex\_} protege el acceso al frame compartido (productor/consumidor).
  \item \texttt{std::atomic<bool> running\_} permite parar el hilo sin data-race.
\end{itemize}

% =============================================================================
\section{Criterio 2 — Calidad del Mapa de Disparidad (25\%)}
% =============================================================================

\subsection*{Lo que pide la rúbrica}
Mapa denso, bordes definidos, bajo ruido. Aplicación justificada de preprocesamiento
(CLAHE) y posprocesamiento (WLS) vinculados a SGBM.

\subsection*{Pipeline completo — \texttt{StereoProcessor.cpp}}

\[
\text{Frame crudo} \xrightarrow{\text{remap}} \text{Rectificado}
\xrightarrow{\text{cvtColor}} \text{Gris}
\xrightarrow{\text{CLAHE}} \text{Contraste mejorado}
\xrightarrow{\text{SGBM}} \text{Disparidad Q16.4}
\xrightarrow{\text{WLS}} \text{Disparidad filtrada}
\xrightarrow{\text{reproject}} \text{Nube 3D}
\]

\begin{lstlisting}[caption={Paso 1: Rectificación (corrección de distorsión)}]
// remap aplica los mapas precalculados de calibracion
// para eliminar distorsion de lente y alinear lineas epipolares
cv::remap(rawLeft,  result.leftRect,  calib_->map1x, calib_->map1y, INTER_LINEAR);
cv::remap(rawRight, result.rightRect, calib_->map2x, calib_->map2y, INTER_LINEAR);
\end{lstlisting}

\begin{lstlisting}[caption={Paso 3: CLAHE — preprocesamiento (líneas 88–92)}]
// clipLimit=2.0: limita la amplificacion maxima del histograma local
// tileSize 8x8: divide la imagen en teselas para ecualizacion local
clahe_ = createCLAHE(cfg.claheClipLimit, Size(8, 8));
clahe_->apply(grayL, grayL);
clahe_->apply(grayR, grayR);
\end{lstlisting}

\begin{critbox}[frametitle={Por qué CLAHE antes de SGBM}]
SGBM busca correspondencias por similitud de intensidad de bloques. En zonas con
iluminación desigual o sombras, el contraste local es bajo y el bloque L no parece
el mismo que en R. CLAHE ecualiza el histograma \textbf{localmente} sin saturar
zonas brillantes, lo que mejora la discriminación en bordes y texturas finas.
\end{critbox}

\begin{lstlisting}[caption={Paso 4: SGBM — parámetros clave (líneas 29–35)}]
sgbm_ = StereoSGBM::create(cfg.minDisp, numDisp, blockSz);
// P1 y P2 controlan la suavidad del mapa de disparidad.
// La relacion P2 = 4*P1 es la recomendacion de Hirschmuller.
sgbm_->setP1(8  * 3 * blockSz * blockSz);
sgbm_->setP2(32 * 3 * blockSz * blockSz);
sgbm_->setUniquenessRatio(cfg.uniquenessRatio);
sgbm_->setMode(StereoSGBM::MODE_SGBM_3WAY);
\end{lstlisting}

\begin{lstlisting}[caption={Paso 5: WLS — posprocesamiento (líneas 101–108)}]
// El matcher derecho calcula disparidad R->L, requerida por WLS
rightMatcher_->compute(grayR, grayL, disp16R);
Mat filtered;
// WLS rellena huecos y suaviza preservando bordes (guided filter)
wlsFilter_->filter(disp16L, result.leftRect, filtered, disp16R);
filtered.convertTo(result.disparity, CV_32F, 1.0 / 16.0);
\end{lstlisting}

\begin{critbox}[frametitle={SGBM vs BM — por qué usamos SGBM}]
BM (Block Matching) busca solo horizontalmente. SGBM (Semi-Global Block Matching)
agrega restricciones de suavidad en 8 direcciones (semiglobal), produciendo mapas
más densos y con menos huecos, especialmente en zonas de textura baja.
\end{critbox}

\begin{lstlisting}[caption={Paso 6: Relleno de zona inválida (propagación horizontal)}]
// SGBM deja ~numDisp columnas negras a la izquierda sin datos.
// Se propaga el primer valor valido de cada fila hacia la izquierda.
for (int y = 0; y < result.disparity.rows; ++y) {
    float* row = result.disparity.ptr<float>(y);
    float fill = 0.0f;
    for (int x = 0; x < result.disparity.cols; ++x) {
        if (row[x] > 0.5f) { fill = row[x]; break; }
    }
    if (fill <= 0.0f) continue;
    for (int x = 0; x < result.disparity.cols; ++x) {
        if (row[x] > 0.5f) break;
        row[x] = fill;
    }
}
\end{lstlisting}

\noindent\textbf{Nota de modificación:} si se pidiera rellenar \textbf{verticalmente} (columna por columna),
bastaría hacer el loop externo sobre \texttt{x} y el interno sobre \texttt{y}, usando
\texttt{at<float>(y,x)} en vez de \texttt{ptr<float>(y)[x]}.

% =============================================================================
\section{Criterio 3 — Precisión y Estabilidad de la Distancia (20\%)}
% =============================================================================

\subsection*{Lo que pide la rúbrica}
Reproyección correcta de la matriz 3D, coordenada Z en unidades reales (cm/mm),
lectura estable con Filtro de Kalman 1D, fluctuación $\leq \pm 3\,\text{cm}$ en entorno estático.

\subsection*{Reproyección 3D}

\begin{lstlisting}[caption={Paso 7: reprojectImageTo3D (línea 131)}]
// Q es la matriz de reproyeccion calculada en la calibracion estereo.
// Convierte el mapa de disparidad en una nube de puntos 3D (X, Y, Z en metros).
reprojectImageTo3D(result.disparity, result.points3D, calib_->Q, true);
\end{lstlisting}

\begin{lstlisting}[caption={Medición robusta: mediana de región 10x10 px}]
float StereoProcessor::medianDepthAt(const Mat& points3D, int cx, int cy, int r) {
    // Recolecta todos los Z validos en un cuadrado de radio r
    for (int y = y0; y <= y1; ++y)
        for (int x = x0; x <= x1; ++x) {
            const float z = points3D.at<Vec3f>(y, x)[2];
            // Solo valores finitos en rango plausible (5 cm - 10 m)
            if (std::isfinite(z) && z > 0.05f && z < 10.0f)
                zVals.push_back(z);
        }
    // nth_element: O(n) - equivalente a mediana sin ordenar todo el vector
    std::nth_element(zVals.begin(), zVals.begin() + mid, zVals.end());
    return zVals[mid];
}
\end{lstlisting}

\subsection*{Filtro de Kalman 1D — \texttt{KalmanDepth.cpp}}

\begin{lstlisting}[caption={Filtro de Kalman escalar (las 3 ecuaciones clave)}]
double KalmanDepth::update(double measurement) {
    // PREDICCION: la incertidumbre crece con el tiempo (ruido de proceso q)
    covariance_ += q_;

    // GANANCIA: balancea confianza en prediccion vs. nueva medicion
    const double gain = covariance_ / (covariance_ + r_);

    // ACTUALIZACION: corrige la estimacion con la nueva medicion
    estimate_   += gain * (measurement - estimate_);
    covariance_  = (1.0 - gain) * covariance_;
    return estimate_;
}
\end{lstlisting}

\begin{center}
\begin{tabular}{lll}
\toprule
\textbf{Variable} & \textbf{Significado} & \textbf{Efecto si sube} \\
\midrule
\texttt{q\_} & Ruido de proceso & Kalman confía más en la medición (responde rápido) \\
\texttt{r\_} & Ruido de medición & Kalman confía más en la predicción (más suave) \\
\texttt{gain} & Ganancia Kalman & Entre 0 (ignorar medición) y 1 (ignorar predicción) \\
\bottomrule
\end{tabular}
\end{center}

\noindent La tecla \textbf{C} reinicia el filtro Kalman central (\texttt{kalmanCenter\_.reset()}).
La tecla \textbf{R} graba las mediciones en CSV con columnas \texttt{frame, elapsed\_ms, z\_raw\_cm, z\_kalman\_cm, fps}.

\subsection*{Validación a 30/60/90 cm}
El script \texttt{python/collect\_measurements.py} automatiza la recolección:
50 muestras por distancia, calcula media, desviación estándar y error absoluto,
y verifica que \texttt{std} $\leq 3.0$ cm (criterio de la rúbrica).

% =============================================================================
\section{Criterio 4 — Efecto de Realidad Aumentada (15\%)}
% =============================================================================

\subsection*{Lo que pide la rúbrica}
Efecto computacional original sobre la transmisión rectificada a color.
El efecto debe reaccionar \textbf{dinámicamente} a la distancia medida (escala, posición
o comportamiento). Diseño exclusivo del grupo.

\subsection*{Descripción del efecto — \texttt{ARRenderer.cpp}}

El efecto \textbf{Energy Ball} tiene tres capas superpuestas, todas moduladas por $Z$:

\subsubsection*{Capa 0: Campo de distorsión por profundidad (\texttt{applyDepthField})}

\begin{lstlisting}[caption={Distorsión de onda reactiva a la distancia}]
// Fuerza inversamente proporcional a Z^3: cae muy rapido con la distancia
const float z3      = zFocusMeters * zFocusMeters * zFocusMeters;
const float strength = std::clamp(0.6f / std::max(z3, 1e-4f), 0.0f, 70.0f);
const float speed    = 10.0f / std::max(0.15f, zFocusMeters);

// Desplazamiento senoidal de cada pixel en primer plano
const float ox = strength * std::sin(y * 0.10f + t * speed);
const float oy = strength * std::cos(x * 0.10f + t * speed * 0.75f);
dst[x] = src.at<cv::Vec3b>(sy, sx);  // lee pixel desplazado
\end{lstlisting}

\noindent\textbf{Comportamiento dinámico:} a 20 cm la onda desplaza hasta 70 px a alta velocidad.
A 80 cm la perturbación baja a 2 px. A más de 80 cm no hay efecto (\texttt{strength < 2.0f}).

\subsubsection*{Capa 1: Tinte radial (\texttt{applyPixelVignette})}

\begin{lstlisting}[caption={Color del tinte reactivo a la distancia}]
// Color cambia segun la distancia (en BGR):
if (depthCm < 40.0f)       { tB=20; tG=60;  tR=255; }  // ROJO   (peligro: muy cerca)
else if (depthCm < 80.0f)  { tB=20; tG=220; tR=80;  }  // VERDE  (zona optima)
else                        { tB=240;tG=160; tR=20;  }  // CYAN   (lejos)
\end{lstlisting}

\subsubsection*{Capa 2: Aura pulsante con destellos orbitales (\texttt{drawEnergyAura})}

\begin{lstlisting}[caption={Pulso y velocidad de destellos reactivos a Z}]
// Pulso a 3 Hz - radio oscila ligeramente
const float pulse  = 0.90f + 0.10f * std::sin(phase * 3.0);
// Mas cerca -> mas destellos y orbita mas rapida
const float proximity = std::clamp(20.0f / depthCm, 0.15f, 1.0f);
const int   sparkCount = 6 + int(8.0f * proximity);  // 6 a 14 destellos
const double spinSpeed = 1.0 + 2.0 * proximity;       // hasta 3x mas rapido
\end{lstlisting}

\begin{center}
\begin{tabular}{lll}
\toprule
\textbf{Distancia} & \textbf{Color} & \textbf{Comportamiento} \\
\midrule
$<$ 40 cm (cerca)   & Rojo  & Onda fuerte, 14 destellos rápidos, tinte intenso \\
40–80 cm (óptimo)   & Verde & Onda moderada, 10 destellos, tinte suave \\
$>$ 80 cm (lejos)   & Cyan  & Onda mínima, 6 destellos lentos \\
\bottomrule
\end{tabular}
\end{center}

% =============================================================================
\section{Criterio 5 — Mini-Evaluación Individual (20\%)}
% =============================================================================

\subsection*{Preguntas probables y respuestas clave}

\textbf{1. ¿Qué es la matriz Q y qué hace \texttt{reprojectImageTo3D}?}

$Q$ es la matriz de reproyección $4\times4$ calculada por \texttt{cv::stereoRectify}.
Transforma disparidad $d$ en coordenadas 3D:
\[
\begin{bmatrix} X \\ Y \\ Z \\ W \end{bmatrix} =
Q \cdot \begin{bmatrix} u \\ v \\ d \\ 1 \end{bmatrix}
\quad\Rightarrow\quad Z = \frac{f \cdot B}{d}
\]
donde $f$ es la longitud focal y $B$ la línea base (distancia entre cámaras).
\texttt{reprojectImageTo3D} aplica esta operación a cada píxel del mapa de disparidad,
produciendo la nube de puntos 3D almacenada en \texttt{result.points3D}.

\vspace{6pt}
\textbf{2. ¿Cómo funcionan los hilos en el sistema?}

Cada \texttt{HttpCameraSource} lanza un \texttt{std::thread} (hilo productor) que descarga
JPEG continuamente con libcurl. El hilo principal (consumidor) llama a \texttt{read()}
que toma el último frame con un \texttt{mutex}. Esta arquitectura productora-consumidora
permite que las descargas HTTP no bloqueen el procesamiento de SGBM.

\vspace{6pt}
\textbf{3. ¿Cómo se sintonizan los parámetros de SGBM?}

\begin{itemize}[leftmargin=*,noitemsep]
  \item \texttt{numDisparities}: rango de búsqueda en píxeles (múltiplo de 16). Mayor = ve más lejos pero más lento.
  \item \texttt{blockSize}: tamaño del bloque de comparación (impar $\geq 3$). Mayor = más suave pero pierde detalle.
  \item \texttt{P1, P2}: penalizaciones de suavidad. $P2 = 4 \cdot P1$ es la relación de Hirschmuller. $P2$ alto → transiciones suaves entre planos.
  \item \texttt{uniquenessRatio}: rechaza coincidencias si la segunda mejor está a menos del X\% de la mejor.
  \item \texttt{speckleWindowSize/Range}: elimina manchas aisladas (ruido de sal y pimienta).
\end{itemize}

En el sistema se pueden ajustar \textbf{en tiempo real} con los trackbars de la ventana de controles sin reiniciar.

\vspace{6pt}
\textbf{4. ¿Qué hace la tecla C?}

Reinicia el filtro Kalman central (\texttt{kalmanCenter\_.reset()}), poniendo
\texttt{initialized\_ = false}. La próxima medición se acepta directamente como
nueva estimación inicial, sin arrastrar el historial anterior.
Útil cuando el objeto cambia bruscamente de posición.

\vspace{6pt}
\textbf{5. ¿Por qué usar la mediana y no el promedio para medir Z?}

La mediana es robusta a valores atípicos (\textit{outliers}).
Los píxeles del borde de un objeto tienen disparidades erróneas por la oclusión.
El promedio sería contaminado por esos valores; la mediana los ignora
siempre que más del 50\% de los píxeles en la región sean válidos.

\vspace{6pt}
\textbf{6. ¿Qué diferencia hay entre Z raw y Z Kalman en el HUD?}

\texttt{Z raw} es la salida directa de \texttt{medianDepthAt} (puede variar frame a frame).
\texttt{Z Kalman} es la estimación filtrada: más estable, reacciona con ligero retraso
pero cumple el requisito de fluctuación $\leq \pm 3\,\text{cm}$ de la rúbrica.

\vspace{6pt}
\textbf{7. ¿Por qué el efecto AR actúa solo en el primer plano?}

En \texttt{applyDepthField} hay un umbral: \texttt{const float fgLimit = zFocusMeters + 0.50f}.
Solo los píxeles con $Z <$ foco $+0.5\,\text{m}$ reciben la distorsión.
Píxeles del fondo (pared, suelo) no se tocan, lo que da sensación de que el efecto
envuelve solo al objeto medido.

\vspace{6pt}
\textbf{8. Si te piden cambiar la dirección del relleno (Paso 6) de horizontal a vertical, ¿qué cambias?}

El loop externo pasa de \texttt{for y} a \texttt{for x}, y el interno de \texttt{for x} a \texttt{for y}.
El acceso pasa de \texttt{ptr<float>(y)[x]} (eficiente, row-major) a \texttt{at<float>(y,x)}
(acceso aleatorio, necesario cuando x es el eje externo).

% =============================================================================
\section*{Resumen de archivos por criterio}
% =============================================================================

\begin{center}
\begin{tabular}{p{4cm}p{5.5cm}p{5cm}}
\toprule
\textbf{Archivo} & \textbf{Responsabilidad} & \textbf{Criterio} \\
\midrule
\texttt{ESP32\_CAM\_Stereo.ino} & Firmware: AEC/AGC/AWB, XCLK, endpoints HTTP & 1 \\
\texttt{CameraSource.cpp/.hpp} & Hilo productor, mutex, lectura no-bloqueante & 1 \\
\texttt{StereoCalibration.cpp} & Carga \texttt{stereo\_calib.yml}, mapas de rectificación, Q & 2, 3 \\
\texttt{StereoProcessor.cpp} & remap → CLAHE → SGBM → WLS → reproject & 2, 3 \\
\texttt{KalmanDepth.cpp} & Filtro Kalman 1D escalar & 3 \\
\texttt{ARRenderer.cpp} & Campo de distorsión, tinte radial, aura pulsante & 4 \\
\texttt{DisplayManager.cpp} & HUD, trackbars, ventanas & 1, 2, 3 \\
\texttt{StereoController.cpp} & Bucle principal MVC, CSV logger, teclas & todos \\
\texttt{collect\_measurements.py} & Validación 30/60/90 cm, CSV de resultados & 3 \\
\texttt{find\_cameras.py} & Escaneo de red para localizar ESP32-CAMs & 1 \\
\bottomrule
\end{tabular}
\end{center}

\end{document}
